% chapters/01/section-03.tex
\section{Peta Jalan Belajar}
\label{sec:peta-jalan}

Bab ini menutup orientasi dengan memperlihatkan seluruh perjalanan yang akan
Anda tempuh: apa yang dipelajari pada tiap tahap, mengapa urutannya demikian,
dan apa yang seharusnya sudah Anda miliki pada akhir setiap tahap.

\subsection{Empat Blok Materi}

Materi mata kuliah ini dikelompokkan menjadi empat blok, sebagaimana
diperlihatkan pada Gambar~\ref{fig:peta-jalan}.

\begin{figure}[htbp]
\centering
% chapters/01/figures/peta-jalan.tikz
% Peta jalan belajar: empat blok materi, dari Python/Pygame ke Unity.
\begin{tikzpicture}[
  blok/.style={
    rectangle, rounded corners=3pt, draw=ubPrimary, thick, fill=ubAlt,
    align=center, inner sep=4pt, minimum height=25mm, text width=26mm,
    font=\footnotesize
  },
  blokunity/.style={blok, draw=ubAccent, fill=white},
  panah/.style={-{Stealth[length=2.6mm]}, very thick, draw=ubAccent}
]

\node[blok] (b1) at (0,0)
  {\textbf{Blok I}\\[2pt]\scriptsize Materi 1--5\\[3pt]\scriptsize \textbf{CPMK1}\\[2pt]\scriptsize \itshape game loop, gerak,\\fisika, collision};
\node[blok, right=8mm of b1] (b2)
  {\textbf{Blok II}\\[2pt]\scriptsize Materi 6--9\\[3pt]\scriptsize \textbf{CPMK2}\\[2pt]\scriptsize \itshape AI NPC, FSM,\\pathfinding A*};
\node[blokunity, right=8mm of b2] (b3)
  {\textbf{Blok III}\\[2pt]\scriptsize Materi 10--11\\[3pt]\scriptsize \textbf{CPMK3}\\[2pt]\scriptsize \itshape Unity, level,\\physics pemain};
\node[blokunity, right=8mm of b3] (b4)
  {\textbf{Blok IV}\\[2pt]\scriptsize Materi 12--14\\[3pt]\scriptsize \textbf{CPMK4}\\[2pt]\scriptsize \itshape aset visual, audio,\\playtesting};

\draw[panah] (b1) -- (b2);
\draw[panah] (b2) -- (b3);
\draw[panah] (b3) -- (b4);

% Keterangan teknologi di bawah blok
\draw[decorate, decoration={brace, amplitude=4pt, raise=3pt}, thick, draw=ubPrimary]
  (b1.south west) -- (b2.south east);
\node[font=\scriptsize\bfseries\color{ubPrimary}, below=9mm of $(b1.south)!0.5!(b2.south)$]
  {Python 3 + Pygame};
\node[font=\scriptsize, below=13mm of $(b1.south)!0.5!(b2.south)$]
  {mekanisme ditulis sendiri};

\draw[decorate, decoration={brace, amplitude=4pt, raise=3pt}, thick, draw=ubAccent]
  (b3.south west) -- (b4.south east);
\node[font=\scriptsize\bfseries\color{ubAccent}, below=9mm of $(b3.south)!0.5!(b4.south)$]
  {Unity};
\node[font=\scriptsize, below=13mm of $(b3.south)!0.5!(b4.south)$]
  {memanfaatkan \itshape engine};
\end{tikzpicture}

\caption{Peta jalan belajar: empat blok materi dari Python/Pygame ke Unity}
\label{fig:peta-jalan}
\end{figure}

Blok I dan Blok II dikerjakan dengan Python dan Pygame, sedangkan Blok III dan
Blok IV dikerjakan dengan Unity. Perbedaan mendasar antara kedua tahap bukan
terletak pada bahasanya---Pygame memakai Python dan Unity memakai C\#---melainkan
pada siapa yang menulis mekanisme permainannya. Pada dua blok pertama, Anda
menulisnya sendiri. Pada dua blok berikutnya, \istilah{engine} yang
menyediakannya dan Anda mengarahkannya.

\begin{table}[htbp]
\centering
\caption{Rincian empat blok materi, capaian, dan perkakasnya}
\label{tab:blok-materi-buku}
{\setlength{\tabcolsep}{4pt}\footnotesize%
\begin{tabularx}{\textwidth}{@{}C{1.6cm} C{2.3cm} C{1.6cm} L{4.3cm} Y@{}}
\toprule
\headrow \textbf{Blok} & \textbf{Materi} & \textbf{CPMK} & \textbf{Pokok bahasan} & \textbf{Perkakas} \\
\midrule
Blok I & Materi 1--5 & \CPMK{1} &
\istilah{Game loop} dan manajemen waktu; gerak berbasis \istilah{delta time};
fisika sederhana; deteksi dan penanganan tabrakan. &
Python 3, Pygame, editor kode, Git. \\
\addlinespace
Blok II & Materi 6--9 & \CPMK{2} &
Perilaku mengejar; mesin keadaan (\istilah{FSM}); pencarian jalur A*;
penggabungan keduanya. &
Python 3, Pygame, modul \code{heapq}; lembar \istilah{playtest}. \\
\addlinespace
Blok III & Materi 10--11 & \CPMK{3} &
Antarmuka dan konsep kerja Unity; penyusunan level; fisika \istilah{engine};
kontrol pemain. &
Unity (versi LTS), editor Unity, akun Unity. \\
\addlinespace
Blok IV & Materi 12--14 & \CPMK{4} &
Impor dan konfigurasi aset visual; animasi; antarmuka pengguna; integrasi
audio; pengujian prototipe dan presentasi. &
Unity, aset visual dan audio bebas pakai, lembar \istilah{playtest}. \\
\bottomrule
\end{tabularx}}
\end{table}

\subsection{Mengapa Python dan Pygame Didahulukan}

Pertanyaan yang wajar muncul: jika Unity menyediakan hampir semua yang
dibutuhkan, mengapa tidak langsung memakainya? Ada empat alasan yang
mendasarinya, dan keempatnya bertumpu pada satu prinsip yang sama, yaitu
memahami sistem sebelum memanfaatkannya.

\begin{enumerate}
  \item \textbf{Mekanisme permainan menjadi terlihat.} Pada \istilah{engine},
        \istilah{game loop}, integrasi fisika, dan deteksi tabrakan sudah
        tertanam dan tidak tampak. Ketika program berperilaku aneh, sulit
        menentukan apakah penyebabnya ada pada kode Anda, pada konfigurasi,
        atau pada cara \istilah{engine} bekerja. Dengan menulis sendiri
        ketiga mekanisme tersebut pada Bab~IV sampai Bab~VII, Anda memiliki
        gambaran yang jelas tentang apa yang sebenarnya terjadi.
  \item \textbf{Algoritma AI dipahami secara algoritmik.} Unity menyediakan
        pencarian jalur siap pakai melalui \istilah{AI Navigation}. Memanggil
        perkakas itu tidak menuntut pemahaman tentang biaya simpul, heuristik,
        atau sifat kelengkapan pencarian. Pada Bab~X Anda akan menulis A*
        sendiri sehingga pengaruh pilihan heuristik terhadap hasil dan
        kinerja pencarian menjadi kasat mata \cite{santos-pathfinding}.
  \item \textbf{Kesalahan lebih mudah dilacak.} Tidak ada lapisan abstraksi
        antara kode Anda dan perilaku program. Pesan kesalahan Python menunjuk
        baris yang tepat, dan seluruh keadaan program berada dalam kendali
        Anda. Kemampuan melacak kesalahan ini terbentuk lebih cepat ketika
        jumlah lapisan yang terlibat sedikit.
  \item \textbf{Biaya dan kelayakan perkakas.} Python, Pygame, dan editornya
        dapat dipasang tanpa biaya lisensi dan berjalan pada komputer
        sederhana. Tahap pertama mata kuliah karena itu tidak bergantung pada
        ketersediaan perangkat berspesifikasi tinggi \cite{sweigart-invent}.
\end{enumerate}

Konsekuensi lanjutannya bermanfaat pada Blok III dan Blok IV. Ketika Anda
menggunakan \istilah{engine}, Anda mengetahui pekerjaan apa yang sedang
diambil alih olehnya, kapan perlu mengubah pengaturannya, dan kapan perlu
menulis kode sendiri. Pengetahuan itu pula yang menjadi bahan refleksi pada
Bab~XVII ketika Anda membandingkan pengalaman bekerja dengan kedua
pendekatan.

\subsection{Alur Enam Belas Minggu}

Perjalanan satu semester terbagi menjadi tiga tahap. Tahap pertama menempati
minggu pertama sampai minggu kesembilan dan menuntaskan seluruh materi
Python/Pygame, yaitu Materi 1--9 pada RPS. Tahap kedua adalah asesmen tengah
semester pada minggu kesepuluh. Tahap ketiga menempati minggu kesebelas sampai
minggu keenam belas dan menuntaskan materi Unity, yaitu Materi 10--14, beserta
presentasi dan ujian akhir.

Rincian pokok bahasan setiap bab beserta Sub-CPMK dan minggu pelaksanaannya
tidak diulang di sini karena telah disajikan pada tabel peta bab di bagian
\istilah{Capaian Pembelajaran dan Peta Kompetensi} di awal buku ini. Jadwal
tugas dan asesmen pada tiap tahap diuraikan pada
Bagian~\ref{sec:kurikulum-obe}.

\subsection{Artefak yang Terkumpul pada Akhir Semester}

Setiap blok menghasilkan artefak yang berbeda. Menyadari hal ini sejak awal
membantu Anda bekerja secara berkelanjutan, karena nilai portofolio ditentukan
oleh kesinambungan pekerjaan mingguan, bukan oleh usaha menjelang ujian.

\begin{table}[htbp]
\centering
\caption{Artefak setiap blok dan waktu pengumpulannya}
\label{tab:artefak}
{\setlength{\tabcolsep}{4pt}\footnotesize%
\begin{tabularx}{\textwidth}{@{}C{1.5cm} L{4.8cm} C{2.4cm} Y@{}}
\toprule
\headrow \textbf{Blok} & \textbf{Artefak} & \textbf{Dikumpulkan} & \textbf{Dinilai dengan} \\
\midrule
Blok I & Prototipe permainan 2D berbasis Pygame: jendela dan
\istilah{game loop}, karakter yang bergerak, gravitasi dan lompatan, serta
pengumpulan item dengan sistem skor. &
Minggu 1--5 & Rubrik portofolio praktikum dan butir Kuis 1. \\
\addlinespace
Blok II & Mini proyek AI: satu atau lebih NPC dengan mesin keadaan dan
pencarian jalur A* pada peta berhalangan, disertai laporan \istilah{playtest}
dan presentasi. &
Minggu 9 & Rubrik mini proyek AI. \\
\addlinespace
Blok III & Level Unity yang dapat dimainkan: lantai dan rintangan, kontrol
pemain, dan fisika yang berjalan benar. &
Minggu 11--12 & Penilaian level dan rubrik proyek akhir. \\
\addlinespace
Blok IV & Prototipe Unity lengkap: aset visual dan animasi, antarmuka pengguna,
audio yang terhubung pada peristiwa permainan, laporan pengembangan, serta
demo. &
Minggu 15--16 & Rubrik proyek akhir dan uji fungsional. \\
\bottomrule
\end{tabularx}}
\end{table}

Bobot setiap komponen penilaian dan rubrik lengkapnya diuraikan pada
Lampiran~A.
